%-------------------------------------------------------------------------------
% rtl66
%-------------------------------------------------------------------------------
%
% \file        summary.tex
% \library     Documents
% \author      Chris Ahlstrom
% \date        2026-01-06
% \update      2026-01-08
% \version     $Revision$
% \license     $XPC_GPL_LICENSE$
%
%     This document provides LaTeX documentation for the Rtl66 library.
%
%-------------------------------------------------------------------------------

\section{Summary}
\label{sec:summary}

   Now let's take a quick overview of the two paradigms: the original paradigm
   and the masterbus paradigm.
   Here is a UML-like diagram of the main classes.

\begin{figure}[H]
   \centering 
   \includegraphics[scale=0.5]{masterbus.png}
   \caption{Rtl66 Overview}
   \label{fig:masterbus}
\end{figure}

   \index{paradigm!RtMidi}
   In this diagram, the classes in the \texttt{rtl} namespace mostly
   support just the original RtMidi paradigm.

   \index{paradigm!masterbus}
   In this diagram, the classes in the \texttt{rtl} namespace mostly
   The classes in the \texttt{midi} namespace add a master-buss with
   a single API client handle that supports multiple busses (i.e. ports).

   Take a little time to work through the diagram before moving
   on to the next sections.

\subsection{Summary / RtMidi API}
\label{subsec:summary_rtmidi_api}

   The basic operation of an "rtmidi" application is to create an
   \texttt{rtl::rtmidi\_in} or
   \texttt{rtl::rtmidi\_out} object.
   These have access to common MIDI API and I/O functions from the
   \texttt{rtl::rtmidi} base class.
   They also create the appropriate MIDI API object, such as a
   \texttt{rtl::midi\_alsa} object, and then open the associated
   port.

   The \texttt{rtl::rtmidi} functions forward operations to the
   \texttt{rtl::midi\_api} object accessed via a pointer.
   (This linkage is not shown in that diagram.)

   The one exception to the RtMidi paradigm here is the
   \texttt{rtl::rtmidi\_engine} class.
   It is used in the master-buss paradigm.
   It is used by the master buss, which sets up the MIDI client
   handle and a port to
   handle the paradigm described in the next section.

   In any case, the result of the RtMidi paradigm is a number of
   ports enumerated from the system, each with it's own independent
   MIDI client handle. In ALSA, each port calls a MIDI-ALSA-handler
   function, passing the client handle to it via an ALSA data
   structure.
   
\subsection{Summary / Masterbus API}
\label{subsec:summary_masterbus_api}

   The master-buss API is very similar to the API used by
   \textsl{Seq66}. It has implications for both JACK and ALSA,
   as noted below.

   The basic operation of a "masterbus" application is to create a
   \texttt{midi::masterbus}.
   It uses an \texttt{rtl::rtmidi\_engine} object to do the work.
   The \texttt{rtl::rtmidi\_engine} object to do the work.
   It opens the MIDI API, sets up the appropriate I/O callback, if
   applicable.

   The "master" then takes care of enumerating the ports that
   exist in the system, creating
   \texttt{midi::bus\_in} and
   \texttt{midi::bus\_out} objects for each port.

   These "busses" are entered into an input or an output
   \texttt{midi::busarray}, where they can be polled or
   handle sending MIDI events.

   \index{issue!masterbus and ALSA}
   For ALSA, as noted in the previous section,
   each port calls the MIDI-ALSA-handler function.
   However, the ALSA API is \textsl{not} thread-safe.
   Therefore, the masterbus must either process the input
   ports in some serial fashion, or use an input handler protected
   by a mutex.  Currently, serial processing is used.

   \index{issue!masterbus and JACK}
   For JACK, one issue with the master-buss paradigm is that output is
   ragged as the JACK frame size is increased. This issue is something
   we will start testing and see how it can be solved.

%-------------------------------------------------------------------------------
% vim: ts=3 sw=3 et ft=tex
%-------------------------------------------------------------------------------
